\section{Introduction}
\label{sec:intro}


Anxiety and depression are among the most widespread mental health disorders, significantly affecting global public health. According to the World Health Organization (WHO), depression is one of the worldwide leading causes of disability, with over 280 million people affected by it \citep{who2023depression}, while anxiety disorders are the most prevalent mental disorders in the world, impacting over 300 million individuals annually~\citep{who2023anxiety}. These disorders not only reduce individual quality of life but also impose substantial economic burdens on public healthcare systems.

Anxiety and depression have distinct diagnostic criteria. Since they frequently co-occur, thier clinical assessment and treatment can be quite complicated though~\citep{kessler2003epidemiology}. Anxiety is typically characterized by excessive worry and heightened physiological arousal, whereas depression involves feelings of persistent sadness, loss of interest, and cognitive symptoms~\citep{apa2013dsm}. However, shared mechanisms, such as heightened emotional dysregulation and stress susceptibility, suggest an overlapping origin.

One framework that offers insight into anxiety's persistence and overlap with depressive symptoms is the Contrast Avoidance Model (CAM). The CAM posits that individuals with anxiety disorders are motivated to maintain a state of heightened negative affect to avoid sharp emotional shifts when confronted emotionally taxing situations~\citep{newman2011contrast}. This sustained negativity may serve as a maladaptive strategy for reducing unpredictability in emotional experiences but can also perpetuate chronic stress and predispose individuals to depression~\citep{llera2014emotional}.

Empirically validating the CAM is essential to advancing our understanding of anxiety and its relationship with depression. By establishing robust evidence for the mechanisms proposed by the CAM, researchers can better differentiate it from other theoretical frameworks and identify specific pathways for intervention. For instance, interventions targeting the cognitive and emotional strategies associated with contrast avoidance may help disrupt the cycle of sustained worry and emotional distress. Furthermore, empirical support for the CAM could inform personalized treatment approaches, optimizing therapeutic outcomes for individuals with anxiety and co-occurring depressive symptoms.

Attempts at validating the CAM have been made before. For example, \citet{cam_questionnaire} developed a questionnaire that empirically shows the workings of the CAM. Further, \citet{rutter2024} use public Twitter data to investigate differences in negative affect variability and levels between anxiety and depression patients. One advantage of their approach over the questionnaire is that they did not use self-reported data (i.e., responses to a survey) but rather sentiment analysis as a measure of negative affect on public Tweets. 

\citeauthor{rutter2024} found that the anxiety cohort did show lower levels of negative affect variability than the depression cohort, which is in keeping with the CAM. They used tweets data for individuals and applied the Valence Aware Dictionary for sEntiment Reasoning (VADER) \citep{vader} to obtain negative affect scores for each tweet. Then, they grouped the individuals into an anxiety and a depression cohort depending on the availability of a tweet indicating clinical diagnosis of either. They computed an overall mean negative affect level and a negative affect standard deviation for each individual and compared the distributions grouped by cohort membership.

Although their research provided valuable insights into empirical assessment of the CAM, their methods exhibit several weaknesses. First, VADER as a rule- and dictionary-based sentiment analyzer is a relatively simple sentiment analysis tool that does not consider the broader context of the text in which it was written. While VADER does handle tasks like negation to a certain degree, it can often fail with more complex sentences and sarcasm \citep{understandingVADER}. This could prove to be a significant weakness to the methodology.

Second, \citeauthor{rutter2024} computed negative affect levels and variability as a mean and standard deviation across the entire time series for each individual. Thus, they fail to capture short-term fluctuations in negative affect levels.

Third, and possibly most importantly, they did not control for external influential factors that determine negative affect levels. For example, \citet{bathina2021Covid} showed that negative affect levels and overall well-being are heavily influenced by global events such as the COVID-19 pandemic and the \textit{Black Lives Matter} riots in the United States. Further, \citet{tenthij2020circadian} found evidence that negative affect levels as displayed by online activity depend on the time of day. In particular, online activity is related to the circadian rhythm of individuals. \citet{dzogang2017seasonal} showed how mood levels depend on the seasons, with \citet{denissen2008weather} further underlining this by showing how mood depends on weather, which in turn is related to seasonality. These are just a few of the studied predictors of mood levels, which were all disregarded by the previous study. 

Additionally, this study aims to address the limitations of prior research by implementing methodological improvements that offer a more nuanced analysis of negative affect variability in individuals with anxiety and depression. By applying moving window aggregations, we aim to capture short-term fluctuations in emotional states rather than relying solely on static measures like mean and standard deviation over an extended period. Moving window aggregations are powerful tools for extracting localized patterns and trends within sequential data. By summarizing a fixed-length segment of data at each step—typically using functions such as mean, median, or variance—this technique enables researchers to observe how a system's behavior evolves over time without losing resolution across the full dataset. This dynamic approach enables a more precise assessment of affective instability, a key characteristic of both anxiety and depression.  

Furthermore, by incorporating linear mixed-effects models, we seek to account for individual-level variability and control for external influences that may confound the relationship between affect variability and diagnostic group membership. Unlike simple mean comparisons, mixed-effects models allow us to analyze the impact of contextual factors such as time of day, seasonal trends, and major global events on negative affect levels. This approach enhances the robustness of our findings by isolating the unique contribution of contrast avoidance to affective patterns in anxiety and depression.  

Ultimately, this study aims to refine our understanding of the Contrast Avoidance Model and its empirical validity. By addressing key methodological concerns and leveraging improved analytical techniques, we strive to provide more reliable insights into the emotional dynamics of anxiety and depression. These findings have the potential to inform clinical interventions by identifying more precise targets for therapy and improving the personalization of treatment strategies for individuals struggling with anxiety-related and depressive symptoms.

Therefore, the goal of this article is to build upon the findings of \citeauthor{rutter2024} and further test their validity. This is done by investigating two research questions. 

\begin{itemize}
    \item Can the findings of \citeauthor{rutter2024} be replicated by using moving window aggregations of negative affect levels?
    \item Can the findings of \citeauthor{rutter2024} be replicated by using linear mixed effects models to build a predictive model of negative affect levels and investigating the unexplained variance?
\end{itemize}




